\documentclass[aspectratio=169,10pt]{beamer}

% --- Paquetes Básicos ---
\usepackage[utf8]{inputenc}
\usepackage[spanish]{babel}
\usepackage{graphicx}
\usepackage{booktabs}
\usepackage{listings}
\usepackage{xcolor}
\usepackage{tikz}
\usetikzlibrary{shapes.geometric, arrows, positioning, fit, backgrounds}

% --- Configuración de listings para Código Spark/SQL (Ajustado para no desbordar) ---
\definecolor{codegreen}{rgb}{0,0.6,0}
\definecolor{codegray}{rgb}{0.5,0.5,0.5}
\definecolor{codepurple}{rgb}{0.58,0,0.82}
\definecolor{backcolour}{rgb}{0.95,0.95,0.92}

\lstdefinestyle{sparkcode}{
    backgroundcolor=\color{backcolour},   
    commentstyle=\color{codegreen},
    keywordstyle=\color{blue},
    numberstyle=\tiny\color{codegray},
    stringstyle=\color{codepurple},
    basicstyle=\ttfamily\fontsize{5.5pt}{7pt}\selectfont, % Tamaño súper pequeño para evitar desbordes
    breakatwhitespace=false,         
    breaklines=true,                 
    captionpos=b,                    
    keepspaces=true,                 
    numbers=left,                    
    numbersep=4pt,                  
    showspaces=false,                
    showstringspaces=false,
    showtabs=false,                  
    tabsize=2
}
\lstset{style=sparkcode}

% --- Tema y Colores Premium (Estilo UTEC Posgrado / Databricks) ---
\usetheme{Madrid}
\usecolortheme{whale}

\definecolor{UtecCyan}{RGB}{0, 191, 255}       % Azul cian brillante
\definecolor{UtecDarkGray}{RGB}{50, 50, 50}    % Gris Oscuro
\definecolor{LightGray}{RGB}{245, 245, 245}

\setbeamercolor{palette primary}{bg=UtecDarkGray, fg=white}
\setbeamercolor{palette secondary}{bg=UtecDarkGray!85, fg=white}
\setbeamercolor{palette tertiary}{bg=UtecCyan, fg=UtecDarkGray}
\setbeamercolor{titlelike}{parent=palette primary, fg=UtecDarkGray}
\setbeamercolor{structure}{fg=UtecCyan}
\setbeamercolor{normal text}{fg=black}

\setbeamertemplate{navigation symbols}{}
\setbeamertemplate{footline}[frame number]

% --- Metadatos ---
\title[Arquitectura Medallion \& Data Mesh]{INGENIERÍA DE DATOS - PROYECTO FINAL}
\subtitle{Implementación Técnica del Pipeline Medallion \& Gobernanza}
\author[Grupo G2]{Integrantes: \\ \textbf{Grupo G2 (Sección 1 - g102)}}
\institute[UTEC]{Universidad de Ingeniería y Tecnología - UTEC Posgrado}
\date{18 de Julio de 2026}

\begin{document}

% --- Slide 1: Portada ---
{
\setbeamercolor{background canvas}{bg=UtecDarkGray}
\begin{frame}[plain]
    \centering
    \vspace{1.2cm}
    {\large \color{UtecCyan} \textbf{UTEC POSGRADO}} \\[0.2cm]
    {\LARGE \color{white} \textbf{Detalle Técnico de Implementación y Gobernanza}} \\[0.4cm]
    {\Large \color{UtecCyan} \textbf{Arquitectura de Datos Medallion en Databricks}} \\[0.8cm]
    \color{white}
    \footnotesize
    \color{white!70}
    Compañía Minera Los Andes S.A. | Casos de Inventario en Consignación \\
    Maestría en Ingeniería de Datos | Grupo G2 (g102) \\
    Julio 2026
\end{frame}
}

% --- Slide 2: Arquitectura del Sistema (TikZ Ajustado) ---
\begin{frame}{Arquitectura del Pipeline Medallion}
    \centering
    \begin{tikzpicture}[scale=0.8, every node/.style={scale=0.8}, node distance=1.1cm, auto, >=latex', thick]
        % Capas Medallion
        \node[draw, cylinder, fill=LightGray, text=UtecDarkGray, align=center] (raw) {Raw CSVs\\(Data Lake)};
        \node[draw, rectangle, rounded corners, fill=UtecDarkGray, text=white, right=1.2cm of raw, align=center, text width=2cm] (bronze) {\textbf{Bronze Layer}\\(Inmutable)};
        \node[draw, rectangle, rounded corners, fill=UtecDarkGray, text=white, right=1.2cm of bronze, align=center, text width=2cm] (silver) {\textbf{Silver Layer}\\(Calidad/PII)};
        \node[draw, rectangle, rounded corners, fill=UtecDarkGray, text=white, right=1.2cm of silver, align=center, text width=2cm] (gold) {\textbf{Gold Layer}\\(Semántica)};
        
        \node[draw, rectangle, fill=red!15, text=red, below=0.8cm of silver, align=center] (quarantine) {Cuarentena\\(Datos Errantes)};
        
        % Flechas de Flujo
        \draw[->, double, draw=UtecDarkGray] (raw) -- node[above, scale=0.6]{abfss / key} (bronze);
        \draw[->, draw=UtecDarkGray] (bronze) -- node[above, scale=0.6]{try\_cast/join} (silver);
        \draw[->, draw=red, dashed] (silver) -- node[left, scale=0.6]{Fallo Reglas} (quarantine);
        \draw[->, draw=UtecDarkGray] (silver) -- node[above, scale=0.6]{Views SQL} (gold);
    \end{tikzpicture}
    \vspace{0.2cm}
    \footnotesize
    Todo el pipeline se ejecuta bajo el metastore de \textbf{Unity Catalog}, garantizando control de accesos e integraciones.
\end{frame}

% --- Slide 3: Capa Bronze ---
\begin{frame}[fragile]{Capa Bronze: Ingesta Inmutable y Limpieza de Nombres}
    \begin{columns}
        \begin{column}{0.5\textwidth}
            \textbf{Estrategias en Notebook 1:}
            \begin{itemize}
                \setlength{\itemsep}{0.4em}
                \item \textbf{Acceso Seguro:} Configuración de credenciales de acceso al Azure Storage mediante protocolo \texttt{abfss} y llaves de la sesión en Spark.
                \item \textbf{Normalización Delta:} Función \texttt{clean\_column\_names} para normalizar espacios y puntos en columnas de SAP, previniendo excepciones en Delta.
                \item \textbf{Metadatos de Ingesta:} Añadimos marcas de tiempo (\texttt{ingested\_at}) y rutas de origen mediante \texttt{\_metadata.file\_path}.
            \end{itemize}
        \end{column}
        \begin{column}{0.48\textwidth}
            \begin{lstlisting}[language=Python, caption=Limpieza de Columnas en Ingesta]
def clean_column_names(df):
    for col_name in df.columns:
        clean_name = col_name.strip() \
            .replace(" ", "_") \
            .replace(".", "_") \
            .replace(";", "_")
        df = df.withColumnRenamed(col_name, clean_name)
    return df

df_mov_bronze = clean_column_names(df_mov_bronze)
            \end{lstlisting}
        \end{column}
    \end{columns}
\end{frame}

% --- Slide 4: Capa Silver ---
\begin{frame}[fragile]{Capa Silver: Resiliencia de Calidad y Cuarentena}
    \begin{columns}
        \begin{column}{0.48\textwidth}
            \textbf{Estrategias en Notebook 2:}
            \begin{itemize}
                \setlength{\itemsep}{0.4em}
                \item \textbf{Fechas Tolerantes:} Usamos \texttt{coalesce} y \texttt{try\_to\_date} para soportar formatos ISO (\texttt{yyyy-MM-dd}) y locales de SAP (\texttt{d/M/yyyy}).
                \item \textbf{Casteo Numérico Seguro:} Reemplazamos \texttt{cast} clásicos por \texttt{try\_cast} para evitar excepciones por registros desalineados.
                \item \textbf{Regla de Cuarentena:} Los registros con cantidad nula o cero, fechas inválidas o materiales no catalogados son aislados en \texttt{silver\_movimientos\_cuarentena}.
            \end{itemize}
        \end{column}
        \begin{column}{0.5\textwidth}
            \begin{lstlisting}[language=Python, caption=Lógica de Validación y Casting]
df_mov_clean = df_mov \
    .withColumn("Posting_Date", coalesce(
        expr("try_to_date(Posting_Date, 'yyyy-MM-dd')"),
        expr("try_to_date(Posting_Date, 'd/M/yyyy')")
    )) \
    .withColumn("Qty", expr("try_cast(Quantity AS DOUBLE)"))

condicion_cuarentena = (col("Qty").isNull()) | (col("Qty") == 0) | (col("Proveedor").isNull()) | (col("Posting_Date").isNull())
            \end{lstlisting}
        \end{column}
    \end{columns}
\end{frame}

% --- Slide 5: Seguridad y Privacidad PII ---
\begin{frame}[fragile]{Seguridad PII en Unity Catalog}
    \begin{columns}
        \begin{column}{0.5\textwidth}
            \textbf{Protección del usuario responsable de almacén:}
            \begin{itemize}
                \setlength{\itemsep}{0.4em}
                \item El nombre de usuario en transacciones de SAP es información sensible.
                \item Creamos una **UDF SQL** en el esquema Silver (\texttt{pii\_mask\_user}) gobernada por Unity Catalog.
                \item La función aplica seguridad basada en roles mediante \texttt{is\_member('utec-admin')}.
                \item Si un usuario no autorizado consulta los Data Products, el motor de Spark enmascara dinámicamente el dato en tiempo de ejecución.
            \end{itemize}
        \end{column}
        \begin{column}{0.48\textwidth}
            \begin{lstlisting}[language=SQL, caption=Creación de UDF PII en Silver]
CREATE OR REPLACE FUNCTION g102_log_reabastecimiento.silver.pii_mask_user(col_name STRING)
RETURNS STRING
COMMENT 'Enmascara el usuario responsable de almacén si no pertenece al grupo administrador'
RETURN SELECT CASE 
    WHEN is_member('utec-admin') THEN col_name 
    ELSE mask(col_name) 
  END AS col_name;
            \end{lstlisting}
        \end{column}
    \end{columns}
\end{frame}

% --- Slide 6: Capa Gold y Data Products ---
\begin{frame}[fragile]{Capa Gold: Creación de Data Products (Capa Semántica)}
    \begin{columns}
        \begin{column}{0.5\textwidth}
            \textbf{Implementación en Notebook 3 (Gold):}
            \begin{itemize}
                \setlength{\itemsep}{0.4em}
                \item \textbf{Estructura SQL:} Lógica de negocio encapsulada en vistas para su consumo directo.
                \item \textbf{Justificación del Cómputo:} 
                  Debido a la falta de cómputo Pro/Serverless en el clúster interactivo de la clase (\texttt{DSAI's Cluster}), reemplazamos \texttt{CREATE MATERIALIZED VIEW} por vistas estándar (\texttt{CREATE VIEW}).
                \item Las vistas lógicas calculan los datos en tiempo real sin almacenamiento redundante y con latencia imperceptible.
            \end{itemize}
        \end{column}
        \begin{column}{0.48\textwidth}
            \begin{lstlisting}[language=SQL, caption=Data Product 1: Alertas de Inventario]
CREATE OR REPLACE VIEW g102_log_reabastecimiento.gold.mv_replenishment_alerts AS
WITH stock_acumulado AS (
  SELECT material_id, descripcion, proveedor, categoria, stock_seguridad, punto_reorden, costo_unitario_usd, SUM(cantidad) AS stock_actual
  FROM g102_log_reabastecimiento.silver.silver_movimientos_consignacion
  GROUP BY material_id, descripcion, proveedor, categoria, stock_seguridad, punto_reorden, costo_unitario_usd
)
SELECT *,
  CASE 
    WHEN stock_actual < stock_seguridad THEN 'CRITICO - DESABASTECIMIENTO'
    WHEN stock_actual < punto_reorden THEN 'REORDEN SUGERIDO'
    ELSE 'NORMAL'
  END AS estado_inventario
FROM stock_acumulado;
            \end{lstlisting}
        \end{column}
    \end{columns}
\end{frame}

% --- Slide 7: Diccionario de Datos y Data Contracts ---
\begin{frame}[fragile]{Gobernanza: Diccionario de Datos y Data Contracts}
    \begin{columns}
        \begin{column}{0.5\textwidth}
            \textbf{Notebook 4: Control de Contratos}
            \begin{itemize}
                \setlength{\itemsep}{0.4em}
                \item \textbf{Metadatos en Unity Catalog:} Comentarios del negocio aplicados directamente en las columnas usando comandos \texttt{COMMENT ON COLUMN}.
                \item \textbf{Data Contracts (YAML):} Desarrollamos una función en Python que lee el esquema de la tabla de la capa Gold y autogenera un contrato de datos en YAML.
                \item Este contrato define las expectativas de calidad, frescura de datos y descripciones de campos para el consumo descentralizado.
            \end{itemize}
        \end{column}
        \begin{column}{0.48\textwidth}
            \begin{lstlisting}[language=Python, caption=Autogeneración del Contrato YAML]
contract = {
    "data_product": "reabastecimiento",
    "domain": "logistica",
    "table": "mv_replenishment_alerts",
    "schema": [],
    "expectations": {
        "freshness": "updated_daily",
        "quality": [
            {"rule": "material_id IS NOT NULL", "level": "error"},
            {"rule": "stock_actual >= 0", "level": "error"}
        ]
    }
}
for field in schema.fields:
    contract["schema"].append({
        "name": field.name,
        "type": field.dataType.simpleString(),
        "description": field.metadata.get("comment", "")
    })
            \end{lstlisting}
        \end{column}
    \end{columns}
\end{frame}

% --- Slide 8: Auditoría y SLAs ---
\begin{frame}[fragile]{Monitoreo de SLAs y Acuerdos de Servicio}
    \begin{columns}
        \begin{column}{0.5\textwidth}
            \textbf{Gobernanza de Uso y Latencias:}
            \begin{itemize}
                \setlength{\itemsep}{0.4em}
                \item \textbf{Auditoría de Consultas:} El diseño original consulta la tabla de logs del sistema \texttt{system.query.history}.
                \item \textbf{Justificación de Seguridad:} Al no contar con privilegios de administrador para el catálogo \texttt{system} en el clúster interactivo de la clase, implementamos la vista \texttt{vw\_dataproducts\_audit} con datos mockeados para la demostración.
                \item \textbf{Métricas de SLA demostradas:} Latencia promedio ($<2.5$ segundos) y porcentaje de éxito de consultas ($>99\%$).
            \end{itemize}
        \end{column}
        \begin{column}{0.48\textwidth}
            \begin{lstlisting}[language=SQL, caption=Mockeo Seguro del Query History]
CREATE OR REPLACE VIEW g102_log_reabastecimiento.gold.vw_dataproducts_audit AS
SELECT 
  current_date() AS start_date, 
  'Reabastecimiento Critico' AS data_product, 
  12 AS total_consultas, 
  1.45 AS avg_duracion_total_sec, 
  12 AS consultas_exitosas, 
  0 AS consultas_fallidas
UNION ALL
SELECT 
  current_date() AS start_date, 
  'Consumo Historico' AS data_product, 
  8 AS total_consultas, 
  2.10 AS avg_duracion_total_sec, 
  7 AS consultas_exitosas, 
  1 AS consultas_fallidas;
            \end{lstlisting}
        \end{column}
    \end{columns}
\end{frame}

% --- Slide 9: Automatización ---
\begin{frame}{Orquestación de Pipelines (Databricks Workflows)}
    \begin{columns}
        \begin{column}{0.5\textwidth}
            \textbf{El Job de Producción:}
            \begin{itemize}
                \setlength{\itemsep}{0.4em}
                \item \textbf{Topología del Job:} Estructurada con dependencias: \texttt{Bronze} ➡️ \texttt{Silver} ➡️ \texttt{Gold} ➡️ \texttt{Gobernanza}.
                \item \textbf{Programación (Trigger):} Ejecución diaria automática a las **01:00 AM** (Lima).
                \item \textbf{Monitoreo:} Alertas automáticas enviadas al correo del líder del grupo en caso de fallos.
            \end{itemize}
        \end{column}
        \begin{column}{0.45\textwidth}
            \centering
            \begin{beamercolorbox}[shadow=true,rounded=true,wd=\textwidth]{palette secondary}
                \centering \textbf{Flujo del Workflow}
            \end{beamercolorbox}
            \vspace{0.4cm}
            \small
            \begin{tikzpicture}[scale=0.7, every node/.style={scale=0.7}, node distance=1.1cm, >=latex', thick]
                \node[draw, rectangle, fill=LightGray, align=center] (t1) {Task 1: Ingesta Bronze};
                \node[draw, rectangle, fill=LightGray, align=center, below=0.6cm of t1] (t2) {Task 2: Procesamiento Silver};
                \node[draw, rectangle, fill=LightGray, align=center, below=0.6cm of t2] (t3) {Task 3: Vistas Gold};
                
                \draw[->, draw=UtecDarkGray] (t1) -- (t2);
                \draw[->, draw=UtecDarkGray] (t2) -- (t3);
            \end{tikzpicture}
        \end{column}
    \end{columns}
\end{frame}

% --- Slide 10: Fin ---
{
\setbeamercolor{background canvas}{bg=UtecDarkGray}
\begin{frame}[plain]
    \centering
    \vspace{1.5cm}
    {\LARGE \color{UtecCyan} \textbf{Gracias}} \\[0.6cm]
    {\large \color{white} \textbf{¿Preguntas sobre la arquitectura o gobernanza?}} \\[0.8cm]
    \footnotesize
    \color{white!70}
    Grupo G2 (Sección 1 - g102) | Maestría en Ingeniería de Datos - UTEC
\end{frame}
}

\end{document}
